\section{Discussion}
\label{sec:discussion}

\subsection{The single-pass strategy is not subsuming}

Our measurement settles a question the predecessor papers left open: the isolated-pass inferrer's \texttt{InterproceduralAnalyzer}~\cite{edmonds2026inference} does not capture what a compositional refinement would, on any of the five libraries we measured. Two structural shapes are out of reach by
construction: branch-conditional implications that depend on the
caller's control-flow context, and disjunctive obligations from
polymorphic dispatch. The cross-library data shows these are not
edge cases --- branch-conditional implications make up the
\emph{majority} of newly-added preconditions on Lang, IO, Math, and
Vavr, and a substantial quarter on jOOL; disjunctive shapes are
present on every library and dominant where interface inheritance
runs deep (Vavr, 9.8\%).

The implication for the inferrer's downstream consumers
(test generators, IDEs, LLMs) is that running the second pass before
emission delivers a measurably richer spec. Whether the additional
clauses translate into stronger downstream behaviour (more bugs
caught, fewer false positives in static analysis, better test
inputs) is the empirical question for follow-up work; we report only
that the gap exists and that it is structurally explainable.

\subsection{Why the single pass already does so much}

We also note what the single pass \emph{does} cover. The
isolated baseline already produces 11.95\% bytecode overhead under
synthetic specs and 10.07\% under real inferred
specs~\cite{edmonds2026embedding}; a substantial fraction of those
clauses are interprocedural propagations the
\texttt{InterproceduralAnalyzer} produces during isolated inference.
The second pass adds 45{,}740 \emph{new} clauses on top of an
already-rich baseline across the five libraries. The single pass
does the right thing on unguarded call sites and on calls to
single-candidate methods; it is the guarded and polymorphic cases
where the gap appears, and those cases concentrate the additions.

\subsection{Implications for analyser architecture}

For inferrers that target downstream tools where richer specs help (deductive verification with KeY~\cite{key2016book} or OpenJML~\cite{openjml}, test generation with EvoSuite~\cite{fraser2011evosuite} or Randoop~\cite{pacheco2007randoop}, mutation-testing pipelines~\cite{coles2016pitest,papadakis2019mutation}, or LLM-driven consumers~\cite{schafer2023adaptive,lemieux2023codamosa,kang2023libro}), the second pass is clearly worth integrating: an extra 34 seconds total across all five libraries for an average of 34.47\% of methods receiving refined specifications is a strong trade. For inferrers that target downstream tools that consume
specs eagerly (auto-completion, hover text in an IDE), the
runtime cost may not justify the extra clauses; the second pass
is opt-in by design so deployers can choose.

A further architectural implication: an inferrer that wants to keep
its pipeline single-pass can still capture some of the
compositional value by enriching its
\texttt{InterproceduralAnalyzer}'s single-call-site lifting. The
branch-conditional shape, which dominates the gap, requires the
analyser to know the enclosing guard of each call site. That
information is available during the AST walk that the isolated pass
already does; the single-pass analyser could in principle lift its
substituted clause through the guard before adding it. The
polymorphic-dispatch shape is harder to incorporate into a
single-pass analysis because it requires looking up sibling cache
entries at the point of propagation, which is what the second pass
does anyway. Whether the single pass should be extended to handle
the easy half of the gap, or whether the two-pass architecture is
the right factoring, is a design choice the present measurement
makes explicit.

\subsection{Limitations}

The 414-line analyser is a scaffold, not a complete WP transformer.
It does not lift through loops (the inferrer's loop-invariant
analyser handles those in the isolated pass), does not reason about
aliasing (the syntactic substitution treats two distinct names as
distinct values), and does not interact with the SMT verifier that
discharges the final clauses. These are the boundaries of the
present contribution; the comprehensive WP transformer would be a
substantially larger development.

The 665 maximum new preconditions on one Commons Math method is the
worst case we observed across the five-library corpus. The method
in question is a deeply branching numerical fa\c{c}ade with many
guarded calls into validation helpers; its 665 new clauses are
overwhelmingly branch-conditional implications gating different
code paths. Whether the resulting spec is \emph{useful} to a
downstream tool --- i.e.\ whether it is parsimonious enough to be
readable and tight enough to be informative --- is the question
we defer to the downstream-impact study planned as future work.
